\scanpage{421}
\mspage{213}
% Main text, approximate manuscript lines 1--4.
\deleted{Let} Is uniqueness at \unsure{all|all} events clear?
\deleted{\unsure{Choosing a}} Let us consider the extension \inserted{$\mathcal{E}$} of
$\Gamma$ by $L=\mathrm{Ker}((N_a)_{I'}\to\unsure{N_a|N_0})$,
inverse \unsure{image} of the extension $(N_a)_{I'}$
(of $\unsure{N_{a'}|N_{0'}}$ by $L$) via $\Gamma\to\unsure{N_{a'}|N_{0'}}$, we
% Main text, approximate manuscript lines 5--8.
have a \inserted{partial} splitting of this extension \unsure{over}
the subgroup $\Gamma'$ of $\Gamma$, and the assertion is
that this splitting \uline{\unsure{extends}}, \uline{in a way} \uline{that is unique}, to $\Gamma$.
We may suppose $\Gamma'$ normal in $\Gamma$, and we
% Main text, approximate manuscript lines 9--13.
identify $\Gamma'$ with a subgroup of $\mathcal{E}$. \unsure{All}
\deleted{\unsure{sections}} the sections of $\mathcal{E}$ over $\Gamma$ \unsure{are identified}
with \inserted{\unsure{section}} subgroups $\widetilde{\Gamma}$ of $\mathcal{E}$.
For such a $\widetilde{\Gamma}$, we have of course
$\widetilde{\Gamma}\subset\mathrm{Norm}_{\mathcal{E}}(\Gamma')$
moreover \deleted{\unsure{if we} \illegible{2}} \inserted{we obviously have}
% Main text, approximate manuscript line 14.
\[
\mathrm{Norm}_{\mathcal{E}}(\Gamma')\cap L=L^{\Gamma'} \tag{31}
\]
% Main text, approximate manuscript lines 15--21.
and I presume we must have
$L^{\Gamma'}=1$ (which generalizes condition
\textcircled{3} above\ldots) Now \uline{this condition}
\uline{implies uniqueness} -- namely
$\widetilde{\Gamma}=\mathrm{Norm}_{\mathcal{E}}(\Gamma')$ --
and existence \uline{means} that
$\mathrm{Norm}_{\mathcal{E}}(\Gamma')\to\Gamma$
is an epimorphism (hence an isomorphism\ldots).

% Main text, approximate manuscript lines 22--23.
But we will have to try to make precise,
in certain contexts (at least that of
\scanpage{422}
% Main text, approximate manuscript lines 1--4; continuation of scan 421.
arithmetically faithful actions, i.e. $\Gamma\to\boldsymbol{\Gamma}_a$
injective -- which \unsure{corresponds} \unsure{normally|presumably}
to the case of algebraic curves \unsure{defined} over
algebraic extensions of $\mathbb{Q}$ -- the notion
% Main text, approximate manuscript lines 5--9.
of an ``admissible'' action. For this, we
must return to the relation between
(potential) curves and finite coverings,
which also gives a way of
varying $g$.

% Main text, approximate manuscript lines 10--17.
But I feel like first taking up again
\unsure{under|on} another aspect \inserted{(perhaps more general)} the preceding
formalism, which \unsure{can} perhaps also \unsure{be} applied to the case of \uline{discrete} groups
with loops (I am thinking of the formalism
of the $\unsure{U^!|\mathcal{U}^!}$, linked to outer actions of
finite groups on such \unsure{profinite groups} with loops).

% Main text, approximate manuscript lines 18--20.
Let $\mathcal{X}$ be a set \inserted{$\ne\varnothing$}
(\inserted{\unsure{morally}, a} set of the
``points'' of an \deleted{surface} algebraic curve,
or of a topological surface\ldots), \deleted{and $\Gamma$}
\scanpage{423}
\mspage{214}
% Main text, approximate manuscript lines 1--7.
\deleted{a group operating on $\mathcal{X}$.} We suppose
given, for every finite subset $I$ of $\mathcal{X}$,
with complement $U_I=\mathcal{X}\setminus I$, a
\unsure{groupoid} $\Pi_{U_I}$, having $U_I$ as its set
of \unsure{objects}. Moreover, for $J\supset I$ i.e.
$U_J\subset U_I$, we suppose given a homomorphism
of groupoids
% Main text, approximate manuscript line 8.
\[
\Pi_{U_J}\to\Pi_{U_I}
\]
% Main text, approximate manuscript lines 9--14.
which on objects is the inclusion $U_J\unsure{\subset|\hookrightarrow}\unsure{U_J|U_I}$.
We suppose transitivity (strict) i.e. \deleted{\unsure{which}}
a covariant functor \inserted{from the category} of subsets $U$ of
$\mathcal{X}$ complementary to finite subsets \inserted{(with the inclusions)}, to
\unsure{the} category of groupoids, which on \unsure{the set}
of objects coincides with the obvious functor\ldots

% Left marginal note, running beside approximate lines 10--22.
\margin{If $\mathcal{X}$ is finite and $I=\mathcal{X}$, we
\unsure{suppose} \inserted{\uuline{\unsure{all the same|however}}}
$\Pi_{U_I}\;(=\Pi_{\varnothing})$ \unsure{is}
not the empty groupoid, but a groupoid
\deleted{of any kind} (which \unsure{maps} into the preceding ones\ldots)
but we will not be able to require
\unsure{the} strict transitivity.
We \unsure{will need|would need} \inserted{\unsure{perhaps}} \illegible{1} anabelian \unsure{hypotheses}
on $\Pi_{U_I}$ in this case\ldots}

% Main text, approximate manuscript lines 15--18.
We suppose moreover that for every \deleted{$I\in\mathcal{P}_f(\mathcal{X})$}
$i\in\mathcal{X}$, we give ourselves a \deleted{\unsure{groupoid} \illegible{1}}
$\Pi_i$ (or $\Pi_{D_i^*}$),
and \deleted{\unsure{for every}} for $i\in I\in\mathcal{P}_f(\mathcal{X})$,
a homomorphism of groupoids
% Main text, approximate manuscript line 19.
\[
\Pi_{D_i^*}\to\Pi_{U_I}
\]
% Main text, approximate manuscript lines 20--23.
compatible with the transition morphisms
\inserted{$\Pi_{U_I}\to\Pi_{U_J}$ $(I\supset J)$}.
We suppose that for fixed $I$, we find
on $\Pi_{U_I}$ a structure of groupoid
with loops by
\[
\Pi_{D_I^*}=\coprod_{i\in I}\Pi_{D_i^*}\to\Pi_{U_I}.
\]
\scanpage{424}
% Main-text lines 1--2; insertion above line 2.
without moreover excluding the case where $I=\varnothing$, and where \inserted{($\Pi_{\mathcal{X}}$ \unsure{being} connected, say)} \unsure{its|the} genus \uline{is $0$} --- \deleted{\unsure{thus}}
\margin{\unsure{We suppose} \illegible{4} \unsure{that} \illegible{1} $\mathcal{X}$ \illegible{3} \deleted{\unsure{for} $\mathcal{X}$ \illegible{2}} \deleted{\unsure{two points}}}

% Main-text lines 3--8.
We suppose moreover that if $J\supset I$ i.e. $U_J\subset U_I$, \inserted{and for $i\in J\setminus I$,} the composite homomorphism
\[
\Pi_{D_i^*}\longrightarrow\Pi_{U_J}\longrightarrow\Pi_{U_I}
\]
is ``constant'', in such a way that $\Pi_{U_I}$ is \unsure{obtained} from $\Pi_{U_J}$ by ``plugging the holes'' at $J\setminus I$.

% Main-text lines 9--21; the middle of the example is heavily revised.
Up to now, \unsure{all|all} this could be visualized for example by starting from an orientable topological surface \deleted{of the form \unsure{$X\setminus S$}; \illegible{4} of type} $X$, and a subset \inserted{\illegible{1}} $\mathcal{X}$ of $X$ \deleted{\illegible{1}}\inserted{\unsure{meeting}} \deleted{\unsure{every connected component}}\inserted{\unsure{every} connected component}; by taking for $\Pi_I$ the restriction to $\mathcal{X}\setminus I$ of the fundamental groupoid of $X\setminus I$ (if $I\ne\mathcal{X}$; if $\mathcal{X}=I$, which requires $\mathcal{X}$ finite, we will take the fundamental groupoid of $X\setminus I=X\setminus\{\mathcal{X}\}$; which we will have trouble mapping into the others with \uline{strict} transitivity, no matter!

% Main-text lines 22--23; sentence continues on the next scan.
But, we are mainly interested in the case $\mathcal{X}$ infinite. Where we \unsure{take} $X$ an algebraic curve over
\scanpage{425}
\mspage{215}
% Main-text lines 1--3; continuation from scan 424.
an algebraically closed field, $\mathcal{X}$ a subset of $X$ meeting \unsure{every} connected component, and we define the $\Pi_{U_I}$ as before.

% Main-text lines 4--9, with an insertion between lines 6 and 7.
Let now $\Gamma$ be a group which operates on the situation (\uline{strictly}, be it understood) thus on \unsure{the set} $\mathcal{X}$, \deleted{and} the groupoids \unsure{$\Pi_{U_I}$}, \inserted{the $\Pi_{D_i^*}$, (N.B. which it \unsure{permutes} among themselves),} commuting with the homomorphisms \unsure{d|of}
\[
\Pi_{U_I}\longrightarrow\Pi_{U_J}
\]
and the $\Pi_{D_i^*}\longrightarrow\Pi_{U_I}$.
\margin{\deleted{(which it \unsure{permutes} \unsure{among themselves})}}

% Main-text lines 10--13; the wording after Nous is rewritten interlinearly.
(We will have difficulties, if $\mathcal{X}$ is finite, for $\Pi_{U_I}$ when $I=\mathcal{X}$, for the \uline{strict} commutation \deleted{\unsure{of the}} \unsure{$\Pi_{U_{\mathcal{X}}}$}$\longrightarrow\Pi_{U_J}$ \unsure{..} but let us pass on...)
We \deleted{make \unsure{the hypothesis|the hypotheses},} \inserted{suppose defined a notion of \uline{admissible} subgroup of $\Gamma$, (they are in any case $\ne1$) such that \deleted{\illegible{1}} if $\Gamma'$ is admissible \unsure{then} every subgroup $\Gamma''\supset\Gamma'$ also.} \deleted{\unsure{following}} \inserted{We \deleted{\unsure{suppose}}}

% Canceled condition and canceled bracketed development in the lower half.
\deleted{\textcircled{\(1^\circ\)} For $P\in\mathcal{X}$, its stabilizer $\Gamma_P$ in $\Gamma$ is \deleted{$\ne1$}\inserted{admissible} (hence $\ne1$)}
\inserted{\deleted{\unsure{$U_I$ stable under $\Gamma_P$ i.e. $I$ stable}}}

\deleted{\unsure{\textcircled{\(2^\circ\)}} \illegible{1} (Let us note that if $U_I\ni P$, then $P$ is an object of $\Pi_{U_I}$ with an operation of $\Gamma_P$ on it, \deleted{\illegible{1}}\inserted{We} \unsure{find} in this way an \deleted{map}\inserted{operation} of $\Gamma_P$ on \unsure{$\pi_1(U_I,P)$|$\Pi_1(U_I,P)$}, and we suppose}

\deleted{\textcircled{\(2^\circ\)} \unsure{$\pi_1(U_I,P)$|$\Pi_1(U_I,P)$} \deleted{\unsure{call}}\inserted{\unsure{say that a}} subgroup $\Gamma'$ of $\Gamma$ is \textcircled{\(1^\circ\)} If $\Gamma'$ \uline{admissible} if \deleted{\illegible{1}} $\mathcal{X}^{\Gamma'}$ \unsure{meets} \unsure{every} connected component of $\mathcal{X}=U_{\varnothing}$ (which}
\scanpage{426}
% Main-text lines 1--3 continue the canceled draft on scan 425.
\deleted{also means \unsure{simply} that $\mathcal{X}^{\Gamma'}\ne\varnothing$ if $\mathcal{X}$ connected, or again that $\Gamma'$ \unsure{contains} a stabilizer $\Gamma_P$.} We assume

% Main-text line 4 and its small parenthetical addition.
\textcircled{\(1^\circ\)} $\forall P\in\mathcal{X}$, $\Gamma_P\ne\{1\}$, and $\Gamma_P$ of finite index \unsure{in $\Gamma$} \inserted{(i.e. the orbit of $P$ finite)}

% Main-text lines 5--21.
We want conditions which ensure that $\mathcal{X}$ and the $\Pi_{U_I}$ \deleted{\unsure{for}} etc, \unsure{can|may} be reconstructed up to \deleted{\unsure{the help}} canonical \unsure{isomorphism} with the help of the purely groupoid-theoretic ``bare'' or topossic \inserted{(with operators $\Gamma$)} \unsure{corresponding|corresponding} data. \deleted{We \illegible{1}} We can say it in learned topossic language, but we are going to express it in terms of the theory of groups with loops. Let $I$ be a subset of $\mathcal{X}$ stable under \unsure{$\Gamma^*$|$\Gamma'$}, we would like to lay down conditions which \deleted{\unsure{we take} a \unsure{subgroup}} \unsure{allow|will allow} us to express \inserted{\unsure{(for such $I$)}} the situation of the groupoids \inserted{$\Pi_V$} associated with the $V\supset U_I$ (i.e. the $U_J$ with $J\subset I$), \inserted{the $\Pi_{D_i^*}$ $(i\in I)$,} and the operation of $\Gamma$ on them, in terms of the \deleted{\unsure{operation}} sole operation of $\Gamma$ on the outer group with loops \deleted{$\pi_I$ associated} $\pi_1(\Pi_{U_I})$ associated with $I$ ---
\scanpage{427}
\mspage{216}
\margin{\unsure{$\rightsquigarrow$|$\longrightarrow$}}
% Main-text lines 1--5.
or rather (for this is immediate, by the operation of drilling holes), \inserted{in the reverse direction,} show how,
\deleted{\unsure{for a} $J\subset I$ \unsure{invariant} under $\Gamma$}
\deleted{\illegible{5}}
\deleted{\unsure{reconstruct} \illegible{2},}
% Main-text lines 6--11.
subject to anabelianity \unsure{$\zeta$|$\xi$} say of \unsure{$\Pi_{\mathcal X}$|$B_{\Pi_{\mathcal X}}$} (to fix ideas), we can $\pm$ \unsure{reconstruct} $\Pi_{U_I}$ and the action of $\Gamma$ on it \unsure{....|...},
starting from \unsure{$\Pi_{\mathcal X}$} \inserted{(or rather, \deleted{from}\inserted{from the \unsure{topos}} \unsure{$B_{\Pi_{\mathcal X}}$})} and from the action of $\Gamma$ on it.
We want at least an \inserted{\unsure{intrinsic}} description of the \unsure{elements} of \unsure{$\mathcal X$}, \inserted{and of the action of $\Gamma$ on it,} via
% Main-text lines 12--15.
this ``soft'' action --- which, for $X$ connected, again amounts to an outer action of $\Gamma$ on a group with loops (\uline{without} loops!) $\pi$...

% Main-text lines 16--21; the parenthesis continues on the next scan.
A first condition, in weak form, \unsure{it is|is} that (assuming \unsure{$\mathcal X$} connected, and calling : $\pi(I)$ the outer group $\pi_1(\Pi_{U_I})$, for \unsure{every} finite subset $I$ of \unsure{$\mathcal X$}) that \unsure{for} $\pi(I)$ anabelian (which will be the case for $I$ ``large enough'',
\scanpage{428}
% Main-text lines 1--3.
at least if \unsure{$\mathcal X$} is infinite \unsure{so|of} that we can take $I$ of arbitrarily large cardinality...) \deleted{we have} the obvious map
% Deleted display and replacement display, approximately lines 4--5.
\[
\deleted{\unsure{\mathrm{Pt}(B_{\pi(I),\Gamma\natural})|\mathrm{Pt}(B_{\pi(I)},\Gamma\natural)}\longrightarrow\text{classes of}}
\]
\[
2^\circ)\quad
\deleted{\unsure{\mathcal X}}U_I\longrightarrow
\unsure{\mathrm{Pt}(B_{\pi(I),\Gamma\natural})|\mathrm{Pt}(B_{\pi(I)},\Gamma\natural)}
\simeq
\begin{gathered}
\unsure{\inserted{\text{classes of }\pi(I)\text{-conjugation }\unsure{\text{of}|\text{of the}}}}\\
\unsure{\text{germs}}\text{ of}\\
\deleted{\text{of}}\text{ splittings of}\\
\text{the extension }\mathcal E(I)\\
\text{of }\deleted{\unsure{\mathcal E}}\Gamma\text{ by }\pi(I)
\end{gathered}
\]
% Main-text lines 6--15.
\unsure{(compatible|compatible} with the actions of $\Gamma$) be \uline{injective}. But to go further, we would have to be able to give a characterization of the image. Let us note (placing ourselves henceforth in the profinite context) that if we give ourselves on the \unsure{$\Pi(I)$|$\pi(I)$} some \unsure{arithmetizations}, \unsure{compatible} with the outer homomorphisms of plugging holes $\pi(I)\longrightarrow\pi(J)$, and with the action of $\Gamma$ on the $\pi(I)$, we therefore have a
% Main-text lines 16--20.
canonical homomorphism from $\Gamma$ into the common group \deleted{of the} $\boldsymbol{\Gamma}$ of \unsure{arithmetically} \unsure{(}outer automorphisms of these \unsure{arithmetizations}.
(and even $\Gamma\longrightarrow$ \unsure{$M_{\pi(\varnothing)=\pi}$}, --- which is \inserted{already} a markedly stronger \inserted{datum}). We can therefore
\scanpage{429}
\mspage{217}
% Main text, approximate lines 1--5.
suppose \unsure{given} a notion of \uuline{admissible}
\uline{lifting} of such an outer arithmetic
action, in such a way that the
preceding \unsure{morphism} is a bijection
\[
 U_I=\unsure{\mathcal X|\mathfrak X}\setminus I
 \xrightarrow{\sim}
 \mathrm{Pt}_{\mathrm{ad}}\bigl(B_{\pi(I)},\Gamma\natural\bigr)
\]
This means moreover, in practice,
that the notion of admissibility satisfies
conditions $1^\circ)$ : $6^\circ)$ \unsure{seen} above.

% Main text, approximate lines 9--12.
As for knowing \unsure{what} should be
called \unsure{operations} \deleted{\unsure{outer}} ``\unsure{admissible|admissible}'' of
$\Gamma$ on an \inserted{outer} \unsure{(}profinite\unsure{)|,} group with loops,
this remains conjectural for the moment.

% Main text, approximate lines 13--21; interlinear revisions retained.
We could by way of experiment
\unsure{conjecture} \deleted{that \unsure{every} \illegible{2}}
\inserted{that the following notion would work}
\unsure{Let us call|Let us call}
``admissible'' every such action
\[
 \Gamma\longrightarrow\widehat{\widehat{\mathcal T}}(\pi)
\]
which respects an arithmetization $a$,
i.e. which factors through the corresponding $N_a$, \unsure{and} such that the image of
$\Gamma\longrightarrow N_a/\widehat{\mathcal T}_a=\boldsymbol{\Gamma}_a$
is \unsure{open|open},
\deleted{\inserted{and that $\pi^{\Gamma'}=(1)$}}
\deleted{\unsure{the conditions}}, \inserted{and that}
\scanpage{430}
% Main text, approximate lines 1--5.
\deleted{\unsure{considered} above \unsure{$1^\circ:6^\circ$} \unsure{are} \unsure{satisfied|satisfied}}

$\forall$ \unsure{homomorphism} of plugging holes at $i\in I(\pi)$,
$\pi\longrightarrow\pi_i'$, $N(\pi)_i\longrightarrow M(\pi_i')$, the \unsure{corresponding} homomorphism
$\Gamma_i$, nor of any open subgroup $\Gamma'$
of $\Gamma_i$, does \unsure{not} normalize an $L_j'$ $(j\in I(\pi_i')=I\setminus\{i\})$.

% Upper left margin; several long strokes run through the whole note.
\margin{\deleted{
this condition is
perhaps \illegible{2} and \unsure{premature}
--- it could be
that for \unsure{every pair|all pairs}
of distinct indices $i,j$ of $I=I(\pi)$,
\inserted{and for the \unsure{corresponding} $L_i,L_j$},
the subgroup \deleted{\illegible{1}} \unsure{$K$} of
$M(\pi)$ \unsure{formed} of the
$g\in M(\pi)$ which \uline{normalize}
$L_i$, and which moreover normalize
$L_j$ \uline{mod} $\mathrm{Ker}(\unsure{\pi\longrightarrow\pi_i'})$,
is ``\unsure{too small}'' \deleted{\unsure{for}} \inserted{\unsure{i.e.}}
that its image in
$\unsure{\boldsymbol{\Gamma}|\Pi}=M(\pi)/\unsure{\widehat{\mathcal T}(\pi)|\widehat{\mathcal G}(\pi)}$
\uline{\unsure{is not}}
of finite index
}}

% Main text, approximate lines 6--14.
\unsure{perhaps} we would even have to impose
$\pi^{\unsure{(\Gamma^{\prime+})|\langle\Gamma^{\prime+}\rangle}}=(1)$, where $\Gamma^{\prime+}$ is the kernel of
$\Gamma'\xrightarrow{\chi}\widehat{\mathbb Z}^{*}$ --- in any case we expect
\unsure{---} that we then have $\pi^{\Gamma'}=(1)$,
and if this were not the case, we would have to
introduce it into the definition --- for
each operation of plugging holes
\unsure{relative} to $I'\subset I(\pi)$, and to the choice of an
$i\in I'$, allowing us to define $\Gamma_{(I',i)}\longrightarrow M(\pi(I'))$.

% Lower left marginal note.
\margin{
More generally,
for $i\in I'\subset I=I(\pi)$,
\deleted{\unsure{it is necessary}} \inserted{\unsure{we impose}} that the
true action of
\unsure{$\Gamma_{I',i}$} on \unsure{$\pi_i(I')$|$\Pi_i(I')$}
\inserted{\unsure{corresponding}} \unsure{to $i$}
normalize no
subgroup $L_j$
$(j\in I\setminus I')$ --- we
are reduced for
this to the case where
$I'=I\setminus\{j\}$, hence
where \unsure{$\pi(I')$|$\Pi(I')$} has \unsure{(only a} single \unsure{class} loops\dots\
(\unsure{if|and} $g\geq1$)
}

% Main text, approximate lines 15--21.
Let us note that a true action of $\Gamma$ on $\pi$
(not \unsure{only|only} \unsure{outer}) which lifts a
given\unsure{$+$} admissible action, is \inserted{\unsure{it too}} admissible
(in the sense that it defines an
\deleted{outer} \uline{admissible} action on a $\pi'$ of
type $(g,\nu+1)$), iff this action \inserted{\unsure{(rather its germ)}} normalizes
no $L_i$ $(i\in I(\pi))$, and that
